\section*{Chapter \ref{ch:iv:options-and-derivatives} --- Options and Derivatives}

\begin{solution}{iq:iv:options-and-derivatives:1}
Parity requires $C - P = 104 - 100e^{-0.01} \approx 4.995$; the quotes give 4.70, so the call is 0.295 cheap relative to the
put. Buy the call, sell the put, short one share, lend $99.005$: the payoff at expiry is zero in every state, and the
position generates about 0.295 today. Check first for a dividend or a borrow cost the parity above ignores (a dividend of
present value 0.30 would explain it), whether the options are American, and whether both quotes are live.
\iqlookfor{parity computed in seconds, the replicating trade, and the reasons it might not be an arbitrage.}
\end{solution}

\begin{solution}{iq:iv:options-and-derivatives:2}
The lower bound is $S - Ke^{-rT} = 50 - 40e^{-0.01} \approx 10.40$, above the quote of 9.50. Buy the call, short the stock,
lend $40e^{-0.01}$: at expiry you hold $\max(S_T - 40, 0) - S_T + 40 \ge 0$, having received about 0.90 today.
\iqlookfor{the lower bound and the trade that exploits it.}
\end{solution}

\begin{solution}{iq:iv:options-and-derivatives:3}
Delta negative, gamma positive, vega positive, theta usually negative. A deep in-the-money European put with positive
rates can have positive theta: it is worth about $Ke^{-rT} - S$, which rises towards $K - S$ as time passes, because the
strike is received later than an American holder could take it.
\iqlookfor{the signs, and the interest-rate effect behind the exception.}
\end{solution}

\begin{solution}{iq:iv:options-and-derivatives:4}
$0.4 \times 100 \times 0.2 \times \sqrt{0.25} = 4.0$; Black--Scholes gives 3.99.
\iqlookfor{the at-the-money approximation from $\varphi(0) \approx 0.4$.}
\end{solution}

\begin{solution}{iq:iv:options-and-derivatives:5}
The 90/100/110 butterfly costs $12 - 2 \times 7 + 1.50 = -0.50$: you are paid to take a position whose payoff is never
negative. Buy it (buy the 90 and 110 calls, sell two 100 calls). Convexity in the strike is violated; the other checks
(prices falling with the strike, slopes above $-1$) pass.
\iqlookfor{the butterfly check and the trade.}
\end{solution}

\begin{solution}{iq:iv:options-and-derivatives:6}
An American call is worth at least the European call, which is at least $S - Ke^{-rT} > S - K$ with positive rates:
exercising early throws away time value and interest on the strike. A put can be worth exercising: at $K = 100$, $S = 10$,
$r = 5\%$, one year, exercising now gives 90, while the European put is worth at most $100e^{-0.05} - 10 \approx 85.12$;
waiting cannot gain more than 10 and delays receiving the strike.
\iqlookfor{the interest argument for calls and a concrete early-exercise case for puts.}
\end{solution}

\begin{solution}{iq:iv:options-and-derivatives:7}
At the money $\Gamma \approx \varphi(0)/(S\sigma\sqrt T)$: about 0.144 for a week and 0.0198 for a year, a factor of about
7.2, which is $\sqrt{52}$. Near expiry the delta swings from 0 to 1 over a narrow range of spot, so gamma concentrates at
the strike (\cref{fig:iv:options-and-derivatives:gamma}).
\iqlookfor{the $1/\sqrt T$ scaling and the picture behind it.}
\end{solution}

\begin{solution}{iq:iv:options-and-derivatives:8}
Gamma P\&L $\tfrac12 \times 50 \times 2^2 = 100$; theta $-\tfrac12 \times 50 \times 100^2 \times 0.04/252 \approx -39.7$;
net about $+60.3$. Break-even when $\tfrac12\Gamma(\Delta S)^2$ equals the theta: $|\Delta S| = S\sigma/\sqrt{252} \approx
1.26$, a one-standard-deviation daily move at the implied volatility.
\iqlookfor{the gamma--theta trade-off and the break-even move as the implied daily move.}
\end{solution}

\begin{solution}{iq:iv:options-and-derivatives:9}
Out of the money, more volatility raises the chance of finishing above the strike: positive vega. In the money, more
volatility raises the chance of falling below: negative vega. The sign changes near the forward; the chapter's code confirms
both signs by finite differences.
\iqlookfor{reasoning from the probability of finishing in the money.}
\end{solution}

\begin{solution}{iq:iv:options-and-derivatives:10}
The digital is $-\partial C/\partial K = \Phi(d_2) - \mathcal V\,\partial\sigma/\partial K \approx 0.460 + 39.7 \times 0.001
\approx 0.500$. The skew lowers the volatility of the calls just above the strike relative to those just below, so the call
spread that replicates the digital costs more than flat volatility implies; pricing a digital off at-the-money volatility
alone would undersell it by about 4 cents on the dollar. A finite difference of calls on the skewed surface gives the same.
\iqlookfor{the digital as a call spread, and the vega-times-skew correction.}
\end{solution}

\begin{solution}{iq:iv:options-and-derivatives:11}
The butterfly $C(95) - 2C(100) + C(105) = 0.8$, divided by the squared spacing 25, gives a density of about 0.032 per unit
around 100, so the probability of ending in the 5-wide bucket around 100 is about $0.032 \times 5 = 0.16$. (A butterfly of
width 5 pays up to 5 at the middle strike; dividing its price by 5 gives the bucket's probability.)
\iqlookfor{Breeden--Litzenberger with finite differences.}
\end{solution}

\begin{solution}{iq:iv:options-and-derivatives:12}
A variance swap is replicated by out-of-the-money options across all strikes, weighted by $1/K^2$, which weights low strikes (puts)
most; with a negative skew those puts carry higher volatility than at the money, so the fair variance exceeds the
at-the-money variance. Integrating the replication formula over the chapter's smile gives a fair volatility of about
20.5\%, against 20.0\% for a flat smile at the same at-the-money level.
\iqlookfor{the $1/K^2$ replication weights and a numerical estimate.}
\end{solution}

\begin{solution}{iq:iv:options-and-derivatives:13}
$C - P = S - D - K = 100 - 2 - 100 = -2$: the put is worth 2 more than the call. With American options parity becomes a pair
of inequalities, $S - D - K \le C - P \le S - K$ with zero rates; the put may be exercised early after the dividend is paid,
and the call just before the dividend if it is deep enough in the money.
\iqlookfor{parity with dividends and its American version.}
\end{solution}

\begin{solution}{iq:iv:options-and-derivatives:14}
The straddle is worth about $0.8\,S\sigma\sqrt{T} = 0.8 \times 100 \times 0.2/\sqrt{252} \approx 1.01$, which you received; you
break even if the stock ends within about $\pm 1.01$. A 3\% move costs $3 - 1.01 \approx 1.99$. The Greeks miss pin risk (not
knowing whether you will be assigned when the stock closes near the strike), jumps into the close, and the gap between the
close and the moment of exercise decisions.
\iqlookfor{the straddle approximation, the break-even, and the expiry risks outside the Greeks.}
\end{solution}
